\begin{explanationbox}
	\textbf{\textcolor{CExplanation}{一句话直觉}}

	向量 $\vec{a}$ 在 $\vec{b}$ 方向上的投影数量，就是 $\vec{a}$ 沿 $\vec{b}$ 方向被“压扁”后得到的有向长度，可正、可负、也可为 $0$；投影向量则是把这个有向长度放到 $\vec{b}$ 方向上的实际向量。

	\medskip
	\textbf{\textcolor{CExplanation}{核心拆解}}
	\begin{description}[style=nextline, leftmargin=2.8em, labelsep=0.5em, itemsep=0.45em]
		\item[\textbf{要点一（投影数量）：}] 投影数量是带符号的标量，由点积除以 $|\vec{b}|$ 得到，体现 $\vec{a}$ 在 $\vec{b}$ 方向上的贡献大小：
		      \[
			      \text{投影数量}=\frac{\vec{a}\cdot\vec{b}}{|\vec{b}|}.
		      \]
		      若 $\vec{a}\ne\vec{0}$，设 $\theta=\angle(\vec{a},\vec{b})$，则也可写成 $|\vec{a}|\cos\theta$。

		\item[\textbf{要点二（投影向量）：}] 投影向量是投影数量乘以 $\vec{b}$ 方向的单位向量，故与 $\vec{b}$ 共线：
		      \[
			      \text{投影向量}
			      =
			      \frac{\vec{a}\cdot\vec{b}}{|\vec{b}|^2}\vec{b}.
		      \]
		      当 $\vec{a}\cdot\vec{b}>0$ 时，投影向量与 $\vec{b}$ 同向；
		      当 $\vec{a}\cdot\vec{b}<0$ 时，投影向量与 $\vec{b}$ 反向；
		      当 $\vec{a}\cdot\vec{b}=0$ 时，投影向量为 $\vec{0}$，没有确定方向。
	\end{description}

	\medskip
	\textbf{\textcolor{CExplanation}{几何本质（有向长度）}}

	投影数量不是普通意义上的“长度”，而是 $\vec{a}$ 在 $\vec{b}$ 方向上的有向长度：锐角时为正，钝角时为负，垂直时为 $0$。投影向量则是将这个有向长度“贴到” $\vec{b}$ 方向上。

	\medskip
	\textbf{\textcolor{CExplanation}{代数意义（点积比例）}}

	点积 $\vec{a}\cdot\vec{b}$ 本身包含投影信息。因为
	\[
		\vec{a}\cdot\vec{b}=|\vec{a}||\vec{b}|\cos\theta,
	\]
	所以除以 $|\vec{b}|$ 就得到投影数量；再除以一次 $|\vec{b}|$，就得到投影向量前面的系数。

	\medskip
	\textbf{\textcolor{CExplanation}{考点价值（常见题型）}}
	\begin{itemize}[leftmargin=*, itemsep=0.5em, topsep=0.4em]
		\item \textbf{考法一（求投影数量）：} 已知向量的坐标或模长与夹角，直接套投影数量公式计算，并注意正负号。
		\item \textbf{考法二（求投影向量）：} 利用投影向量与 $\vec{b}$ 共线，通过系数 $\frac{\vec{a}\cdot\vec{b}}{|\vec{b}|^2}$ 写出投影向量。
		\item \textbf{考法三（求投影向量模长）：} 投影向量的模长为 $\frac{|\vec{a}\cdot\vec{b}|}{|\vec{b}|}$，一定是非负数。
	\end{itemize}

	\medskip
	\textbf{\textcolor{CExplanation}{顿悟点}}

	投影是将向量点积“翻译”为方向贡献的工具。当题目要求“一个向量在另一个方向上的分量”时，优先想到投影公式。

	\medskip
	\textbf{\textcolor{CExplanation}{使用场景}}
	\begin{itemize}[leftmargin=*, itemsep=0.5em, topsep=0.4em]
		\item \textbf{场景一（力的分解）：} 求一个力在某个方向上的分力向量，可用投影向量；求该分力的大小，则取投影数量的绝对值。
		\item \textbf{场景二（正交分解）：} 将向量分解为与已知向量平行和垂直的两部分，先求投影向量即得平行分量。
	\end{itemize}
\end{explanationbox}